\chapter{变量选择与高维数据}

前几章介绍了几种校正方法，用以估计处理 $A$ 对结局 $Y$ 的因果效应，包括分层与结局回归、标准化与参数 g 公式、逆概率加权，以及 g 估计。这些方法以不同方式实施校正，却都依赖同一个条件：校正变量集 $L$ 必须包含足够的信息，使受处理者（$A=1$）与未受处理者（$A=0$）之间达成条件可交换性；等价地说，必须阻断 $A$ 与 $Y$ 之间的全部后门路径，同时不打开其他带来偏倚的路径。

实践中一个常见的问题是：如何挑选用于校正的变量 $L$？当数据分析的目标是因果推断时，本章给出若干挑选变量的准则。因果推断的变量选择标准与预测并不相同，因此预测分析中广为使用的变量选择程序未必能直接搬用到因果分析中。本章综述因果分析中错误选择变量所带来的问题，并勾勒一些实用的建议。

\section{变量选择的不同目标}

贯穿全书的讨论表明，有效的因果推断通常需要对混杂及其他偏倚做校正。当处理 $A$ 与结局 $Y$ 之间的某个关联度量可以部分或全部由混杂因子 $L$ 解释时，就必须把对这些混杂因子的校正纳入数据分析；否则，该关联度量无法解释为因果效应度量。\mnote{即便结局模型包含了 $A$ 对 $Y$ 之效应的全部混杂因子，各混杂因子与结局之间的关联也不能作因果解释，因为我们并未校正这些混杂因子的混杂因子。}

但如果数据分析的目标纯粹是预测，就不需要对混杂做任何校正。如果我们只想量化戒烟 $A$ 与体重增加 $Y$ 之间的关联，只需比较戒烟者与未戒烟者的平均体重增加，直接从数据中估计该关联即可。更一般地，若要为体重增加建立一个预测模型，我们会希望纳入那些能预测体重增加的协变量（如戒烟、基线体重与年收入）。我们不会追问这些协变量是否为混杂因子，因为这里根本没有其效应可能被混杂的处理变量。在预测模型中，我们并不试图赋予任何参数估计以因果解释，因而也不试图校正混杂——混杂这一概念在此根本不适用。\mnote{提醒：混杂是一个因果概念，当估计目标是关联而非因果效应时并不适用。}

预测/关联模型与因果模型之别，已在第 15.5 节讨论。设想临床研究者用结局回归识别心力衰竭的高危病人。目标是分类，而分类是预测的一种形式。这类预测模型的参数未必具有任何因果解释，模型中所有协变量的地位相同，即不存在处理变量 $A$ 与校正变量 $L$。例如，既往住院史可能被识别为预测未来心力衰竭的有用指标，但没有人会建议为了预防心力衰竭而不再让病人住院。识别预后不良的病人（预测），与识别预防或治疗疾病的最佳行动方案（因果推断），是两回事。

为纯粹预测的目的，研究者希望使用能提高预测能力的变量。多数预测算法都含有所谓的调优参数，其取值必须选定以优化预测精度。例如，lasso 与岭回归都含有一个把回归系数向零收缩的正则化参数，而收缩的适当程度（即正则化参数的大小）需要从数据中自适应地选定。对神经网络而言，调优参数包括网络的深度与宽度。调优参数的选择常用交叉验证（见精要 18.2）。交叉验证也用于在相互竞争的算法之间做选择，因为没有任何一种算法在所有数据集上的预测都优于其他算法。

由于某些选择算法（如深度神经网络）是“黑箱”程序，要解释变量如何被选出、算法为何有效，并不总是容易的事。一种观点认为这无关紧要：对特定总体与特定情境下的纯预测目的而言，只要能改进预测，无论可解释性如何，任何算法都可以用。

另一种观点认为，可解释的算法是必需的，因为医生若不能向自己、也向病人说清治疗决策背后的医学理由，他们在做决策时（尤其面对症状组合不典型的病人时）会感到不安。此外，黑箱算法在部署到新情境时可能表现很差，因为它们很可能依赖训练情境中局部性的、往往非因果的特征，而这些特征在新情境中并不存在。最后，黑箱算法可能把（社会上）歧视性的做法固化进模型——这些做法往往有预测力——因为训练集数据正是在这些做法盛行之时收集的。

因果分析需要不同的考量。与预测分析不同，在因果分析中，若想相信治疗效应估计具有因果解释，就必须审慎地挑选混杂因子。自动变量选择程序或许适用于预测，但一般不能用于因果推断。变量选择算法可能选出一些变量做校正，反而在效应估计中引入偏倚。这种偏倚的产生有若干原因，其中一些已在前文讲过，另一些尚未提及。下一节把它们一并总结。

\section{诱发或放大偏倚的变量}

设想我们拥有无限的计算能力，数据集包含准无限多个个体（数据集的行），以及每个个体身上测得的众多变量（数据集的列），其中包括处理 $A$、结局 $Y$，以及数量适中的离散变量 $X$，其中一些可能是 $A$ 对 $Y$ 之效应的混杂因子。在这种情形下，我们可以不受计算、数值或统计上的限制，对数据集中任意多的变量做校正。于是，若目标仅仅是由 $A$ 与 $X$ 预测 $Y$（在标准最小二乘损失下，见技术注记 18.1），我们只要使用 $A$ 与 $X$ 每个联合层中 $Y$ 的平均值，就能最优地预测 $Y$。

\begin{figure}[htbp]
\centering
\begin{tikzpicture}[causal,shorten >=3pt,shorten <=1.5pt,every node/.style={font=\small},scale=1.0]
  \node (A) at (0,0.9) {$A$};
  \node (L) at (2.0,0.9) {$L$};
  \node (Y) at (4.0,0.9) {$Y$};
  \node (U) at (0,-0.5) {$U$};
  \draw[ed] (A) -- (L);
  \draw[ed] (U) -- (L);
  \draw[ed] (U) -- (Y);
\end{tikzpicture}
\caption{}
\label{fig:18.1}
\end{figure}

\renewcommand{\FPname}{精要（Fine Point）}
\begin{finepoint}[回归模型的变量选择程序]
假设我们要拟合一个以预测为目的的回归模型，但数据库含有的潜在预测变量太多——甚至可能多于个体数——把它们全部纳入模型，或者不可能，或者会得到极不稳定的预测。回归模型中有若干办法处理这一问题，许多教材对这些程序有详细描述，例如 Hastie、Tibshirani 与 Friedman (2009) 以及 Harrell (2015) 的书。下面简要勾勒其中一些现有做法。

一种做法是从可得变量中选出一个子集。寻找最优子集在概念上最简单的办法是：先定下模型中变量的个数，再拟合该变量个数下所有可能的模型组合，最后按某个预先指定的准则（如赤池信息准则）挑出最好的一个。然而，变量极多时这种做法在计算上不可行；而且对有限数据集，它也不保证选出预测误差最小的模型。计算上更高效的变量选择方法有前向选择（从没有变量开始，算法每一步加入使改进最大的变量）、后向剔除（从全部变量开始，每一步删去使改进最小的变量）与逐步选择（前向选择与后向剔除的结合）。当无法进一步改进时，变量选择算法终止，这里的“改进”同样按某个预先指定的准则定义。这些算法易于实现，但另一方面，它们并不探索变量的全部可能子集。

子集选择之外的另一种做法是收缩。其思路是修改估计方法，加入一个“惩罚”，迫使模型参数估计（截距除外）比没有惩罚时更接近零。也就是说，多数参数估计被向零收缩。收缩的结果是方差下降、预测更为稳定。最著名的两种收缩方法是岭回归与 lasso，即“最小绝对收缩与选择算子”，它由 Santosa 与 Symes (1986) 提出，后由 Tibshirani (1996) 重新发现。与岭回归不同，lasso 允许某些参数值恰为零。因此，lasso 既是收缩方法，也是子集选择方法。
\end{finepoint}
\renewcommand{\FPname}{精要}

然而，设想我们要无偏地估计二值处理 $A$ 对结局 $Y$ 的平均因果效应，即 $\E[Y^{a=1}]-\E[Y^{a=0}]$。那么协变量校正的目标，就是利用已测变量 $X$ 所含的信息，尽可能消除混杂。我们可以很容易地通过分层/结局回归、标准化/g 公式、逆概率加权或 g 估计，对全部已测变量 $X$ 做校正。那么，是否有理由只校正 $X$ 的一个子集，而非直接校正所有可得的变量 $X$？答案是肯定的。即便在这个理想情境中，我们也要确保某些变量未被选入校正集，因为对它们做校正会引入偏倚。下面的例子说明：当 $X$ 中的某些变量 $L$ 受 $A$ 的因果影响时，会出现什么情况。\mnote{可收缩性提醒：用分层对协变量做校正时，请记住校正后的关联度量可能与未校正的关联度量不同，即便不存在混杂。见精要 4.3。}

设问题的因果结构如图 18.1 的因果图所示（与图 7.7 相同），其中变量 $L$ 是对撞节点。此处平均因果效应 $\E[Y^{a=1}]-\E[Y^{a=0}]=0$，而由于不存在由 $L$ 引起的混杂，它可由 $\E[Y\given A=1]-\E[Y\given A=0]$ 无偏地估计。现在设想我们试图通过对 $L$ 做校正来估计平均因果效应，即用 g 公式
\[
\sum_l \E[Y\given A=1,L=l]\Pr(L=l)-\sum_l \E[Y\given A=0,L=l]\Pr(L=l).
\]
这一对比不同于 $\E[Y\given A=1]-\E[Y\given A=0]$，从而是有偏的，因为 $L$ 既在给定 $A$ 的条件下与 $Y$ 相关，又与 $A$ 边际相关，故 $\Pr(L=l)\neq\Pr(L=l\given A)$。由于即便处理 $A$ 对结局 $Y$ 的因果效应为零，校正了 $L$ 的 $A$-$Y$ 关联仍预期非零，我们说存在\emph{零效应下的选择偏倚}。当我们校正的变量 $L$ 如图 18.2 的因果图所示是对撞节点 $B$ 的子孙时，预期也会出现同样的偏倚。关于含对撞节点及其子孙的因果结构，可回顾第 8 章的更多例子。

\begin{figure}[htbp]
\centering
\begin{tikzpicture}[causal,shorten >=3pt,shorten <=1.5pt,every node/.style={font=\small}]
  \node (A) at (0,0.9) {$A$};
  \node (B) at (1.9,0.9) {$B$};
  \node (L) at (3.6,0.9) {$L$};
  \node (Y) at (5.3,0.5) {$Y$};
  \node (U) at (0,-0.4) {$U$};
  \draw[ed] (A) -- (B);
  \draw[ed] (B) -- (L);
  \draw[ed] (U) -- (B);
  \draw[ed] (U) -- (Y);
\end{tikzpicture}
\caption{}
\label{fig:18.2}
\end{figure}

\begin{figure}[htbp]
\centering
\begin{tikzpicture}[causal,shorten >=3pt,shorten <=1.5pt,every node/.style={font=\small}]
  \node (A) at (0,0) {$A$};
  \node (Y) at (2.0,0) {$Y$};
  \node (L) at (4.0,0) {$L$};
  \draw[ed] (A) -- (Y);
  \draw[ed] (Y) -- (L);
\end{tikzpicture}
\caption{}
\label{fig:18.3}
\end{figure}

\begin{finepoint}[过拟合与交叉验证]
过拟合是回归模型一切变量选择方法的通病：变量被选来尽可能好地预测数据点，却没有考虑到数据中观察到的某些变异纯属随机。结果，模型对用于估计其参数的那些个体预测得很好，但对未用于估计参数的未来个体预测很差。随机森林、神经网络以及其他机器学习算法等预测算法也有同样的问题。

过拟合问题的一个直接解决办法是把样本分成两部分：训练样本，用来运行预测算法（就回归而言，即估计模型参数）；验证样本，用来评估算法预测的准确度。设样本量为 $n$，我们用 $v$ 个个体作验证集，用 $n-v$ 个个体作训练集。使用 lasso 时，训练样本中的收缩程度可由模型在验证样本中的表现来引导。

把样本分成训练与验证两个子样本，明显的缺点是预测算法只用到个体中的一个子集（例如模型参数只在该子集中估计），这增大了方差。一个解决办法是把分样本的过程重复多次，从而提高预测算法实际使用的个体数。然后可以把算法在各验证样本上的平均表现作为其预测精度。这一程序称为交叉验证或样本外检验。交叉验证有若干不同形式。

一种称为“留 $v$ 交叉验证”的程序，分析把样本分成大小为 $v$ 的验证样本与相应训练样本的所有可能划分。然而，当 $n$ 与 $v$ 中等偏大时，检视所有这些划分在计算上可能不可行。两个可能的补救办法是：(i) 取 $v=1$；(ii) 只评估一部分划分。例如在“$k$ 折交叉验证”中，样本被分成 $k$ 个等大的子样本；然后依次把每个子样本用作验证样本，其余 $k-1$ 个子样本用作其训练样本。常见的选择是 $k=10$。关于交叉验证及相关技术的描述，见 Hastie、Tibshirani 与 Friedman (2009) 的书。

当有海量训练数据时（如语音识别、图像），基于多层神经网络的深度学习算法往往似乎几乎不受过拟合之害（Goodfellow, Bengio, Courville 2016）。深度神经网络是一个参数模型，常含成千、上百万甚至上十亿个未知参数，因而往往能精确拟合训练数据。令人惊讶的是，拟合出的模型仍能以很小的误差成功预测未来个体（抽自同一总体）的结局。试图解释这一现象，是当今机器学习中最活跃的研究领域之一。
\end{finepoint}

当我们校正的变量是受处理影响的非对撞节点时，也可能出现选择偏倚，如图 18.3 因果图中的变量 $L$。此处平均因果效应 $\E[Y^{a=1}]-\E[Y^{a=0}]\neq 0$，同样因不存在由 $L$ 引起的混杂，它可由 $\E[Y\given A=1]-\E[Y\given A=0]$ 无偏地估计。然而，如果我们试图通过对 $L$ 做校正（仿佛它是处理前变量）来估计平均因果效应，g 公式的对比将不同于 $\E[Y\given A=1]-\E[Y\given A=0]$，理由与上一段相同。

现在设想从 $A$ 到 $Y$ 的箭头不存在，即“$A$ 对 $Y$ 无效应”的零假设为真，从而 $\E[Y^{a=1}]-\E[Y^{a=0}]=0$。那么 $A$ 与 $Y$ 独立（无论边际还是在给定 $L$ 的条件下），g 公式的对比为零，因而无偏。这一结果的关键原因在于：在零假设下，$A$ 不再对 $L$ 有因果效应。也就是说，与图 18.1、18.2 不同，在图 18.3 中校正 $L$ 只在 $A$ 对 $Y$ 有非零因果效应时才导致选择偏倚。我们于是说存在\emph{备择假设下}（或\emph{偏离零假设时}）的选择偏倚（见第 6.5 节）。

当校正变量受处理 $A$ 影响、又影响结局 $Y$ 时，我们说该变量是\emph{中介}。考虑图 18.4 的因果图，其中中介 $L$ 位于从处理 $A$ 到结局 $Y$ 的一条因果路径上。\mnote{在图 18.4 中，校正 $L$ 阻断了路径 $A\to L\to Y$，但未阻断路径 $A\to Y$。因此校正了 $L$ 的 $A$-$Y$ 关联是 $A$ 对 $Y$ 总效应的有偏估计量，却是 $A$ 对 $Y$ 未经 $L$ 中介的直接效应的无偏估计量。}校正了中介 $L$（或其子孙）之后的 $A$-$Y$ 关联，将不同于处理 $A$ 对结局 $Y$ 的效应，因为该校正阻断了经由 $L$ 的那部分效应。当所关心的对比是 $A$ 对 $Y$ 的平均因果效应时，这一问题有时被称为对中介的\emph{过度校正}。

\begin{figure}[htbp]
\centering
\begin{tikzpicture}[causal,shorten >=3pt,shorten <=1.5pt,every node/.style={font=\small}]
  \node (A) at (0,0) {$A$};
  \node (L) at (2.0,0) {$L$};
  \node (Y) at (4.0,0) {$Y$};
  \draw[ed] (A) -- (L);
  \draw[ed] (L) -- (Y);
  \draw[ed] (A) to[out=40,in=140] (Y);
\end{tikzpicture}
\caption{}
\label{fig:18.4}
\end{figure}

上述诱发偏倚的变量有一个共同特征：它们受处理影响，因而都是处理后变量。人们或许会想，那就应当始终避免校正处理 $A$ 之后发生的变量。这条“不校正治疗后变量”的规则容易遵守，因为校正变量与处理的时间先后通常已知。遗憾的是，照这条简单规则行事，可能把有用的校正变量排除在外，正如我们在精要 7.4 中所讨论的。考虑图 18.5 的因果图。变量 $L$ 是治疗后变量，却可用来阻断处理 $A$ 与结局 $Y$ 之间的后门路径。因此，校正了 $L$ 的 $A$-$Y$ 关联是 $A$ 对 $Y$ 之效应的无偏估计量，而未校正的 $A$-$Y$ 关联则是有偏估计量。要点在于：因果图并不在意时间顺序。因此，当 $A$ 不影响 $L$ 时，无论 $L$ 在时间上先于还是后于 $A$，正确的分析都必须相同。\mnote{关于运用专家知识做校正的一个例子，见 Hern\'an 等 (2002)。}

\begin{figure}[htbp]
\centering
\begin{tikzpicture}[causal,shorten >=3pt,shorten <=1.5pt,every node/.style={font=\small}]
  \node (A) at (0.6,1.0) {$A$};
  \node (L) at (2.8,1.0) {$L$};
  \node (Y) at (4.4,1.0) {$Y$};
  \node (U) at (0,0) {$U$};
  \draw[ed] (U) -- (A);
  \draw[ed] (U) -- (L);
  \draw[ed] (L) -- (Y);
\end{tikzpicture}
\caption{}
\label{fig:18.5}
\end{figure}

问题在于，即便知道变量的时间顺序，我们仍无法从数据判断 $A$ 是否影响 $L$。事实上，给定时间顺序 $A$、$L$、$Y$，任何不含独立性的 $(A,L,Y)$ 联合分布，都与若干种因果图相容。所以是否校正 $L$ 的决定必须依据数据之外的信息。也就是说，是否校正 $L$ 无法通过任何只依赖统计关联的自动化程序来确定。例如，如第 7 章所述，仅凭数据无法区分对撞节点与混杂因子。相反，把诱发偏倚的变量排除出校正集，需要以关于问题因果结构的专业知识为指导。

下面转向时间上先于处理 $A$ 的变量 $L$ 的校正问题，即此时的时间顺序为 $L$、$A$、$Y$。为简单起见，设样本量极大，远远超过可供校正的协变量 $X$ 的数目。于是任何估计量的方差都可忽略，唯一的问题是偏倚。在这种情形下，人们普遍相信：校正全部可得的处理前协变量的估计量将使偏倚最小。然而，这一信念是错误的，原因有两个，彼此独立。

考虑图 18.6 的因果图（与图 7.4 相同），其中含一个处理前变量 $L$。由于 $L$ 是 $A$ 到 $Y$ 某条路径上的对撞节点，对它做校正会引入选择偏倚，第 7 章称之为 M 偏倚。同样，观测数据无法区分混杂因子与对撞节点，因此只能依靠手头的外部信息来判断是否校正处理前变量 $L$。事实上，$L$ 也可能既是混杂因子又是对撞节点——若图 18.6 中有一条从 $L$ 指向 $A$ 的箭头——这意味着无论是否校正 $L$，平均因果效应都无法识别。

\begin{figure}[htbp]
\centering
\begin{tikzpicture}[causal,shorten >=3pt,shorten <=1.5pt,every node/.style={font=\small}]
  \node (U1) at (0,1.9) {$U_1$};
  \node (L)  at (0,0.9) {$L$};
  \node (U2) at (0,-0.2) {$U_2$};
  \node (A)  at (2.4,0.9) {$A$};
  \node (Y)  at (4.4,0.9) {$Y$};
  \draw[ed] (U1) -- (L);
  \draw[ed] (U1) -- (Y);
  \draw[ed] (U2) -- (L);
  \draw[ed] (U2) -- (A);
  \draw[ed] (A) -- (Y);
\end{tikzpicture}
\caption{}
\label{fig:18.6}
\end{figure}

还有一个理由要求避免不加区分地校正处理前变量：\emph{偏倚放大}，这是本书尚未描述过的一种现象。考虑图 18.7 的因果图（与图 16.1 相同），它表示处理 $A$ 对结局 $Y$ 的因果效应被未测量变量 $U$ 混杂的情形。由于 $U$ 不在数据中，我们无法校正 $U$，混杂难以处理。对变量 $Z$ 做校正（用上面的 g 公式，把 $L$ 换成 $Z$）并不能消除混杂，因为 $Z$ 不在从处理 $A$ 到结局 $Y$ 的任何后门路径上。事实上，$Z$ 是一个工具变量——在第 16 章所述的某些情形下可用于工具变量估计——因而对用 g 公式做直接的混杂校正毫无用处。

\begin{figure}[htbp]
\centering
\begin{tikzpicture}[causal,shorten >=3pt,shorten <=1.5pt,every node/.style={font=\small}]
  \node (Z) at (0,0) {$Z$};
  \node (A) at (2.2,0) {$A$};
  \node (Y) at (4.4,0) {$Y$};
  \node (U) at (0,-1.3) {$U$};
  \draw[ed] (Z) -- (A);
  \draw[ed] (A) -- (Y);
  \draw[ed] (U) -- (A);
  \draw[ed] (U) -- (Y);
\end{tikzpicture}
\caption{}
\label{fig:18.7}
\end{figure}

有趣的是，尽管 $Z$ 不能用来校正掉由 $U$ 所致的混杂偏倚，校正工具变量 $Z$ 却可能放大由 $U$ 所致的混杂偏倚。也就是说，校正了 $Z$ 的 $A$-$Y$ 关联，可能比未校正 $Z$ 的 $A$-$Y$ 关联离 $A$ 对 $Y$ 的效应更远。\mnote{若结构方程模型中对应于该因果图的方程全部为线性，则偏倚放大必然发生（Wooldridge 2010, Pearl 2011）；但在更接近现实的情境中也可能出现（Ding et al. 2017）。}这种因校正工具变量 $Z$ 而产生的偏倚放大，常被称为 $Z$ 偏倚，它正是我们应避免校正像 $Z$ 这类工具变量的理由。不过，偏倚放大并非必然发生：校正 $Z$ 也可能减小由未测量变量 $U$ 所致混杂带来的偏倚。一般而言，无法知道校正工具变量会放大还是减小偏倚。

总之，即便没有计算上的限制、样本量准无限，也不宜校正全部可得的变量 $X$。理想情况下，校正集不应包含任何可能引入或放大偏倚的变量。由于这些诱发偏倚的变量无法由纯统计算法从经验上识别，需要专家知识来指导变量选择。

\section{因果推断与机器学习}

本章余下部分，我们假定已经以某种方式确保 $X$ 中不含可能引入或放大偏倚的变量（即不含一旦校正就会破坏条件可交换性的变量），同时仍包含 $A$ 对 $Y$ 之平均因果效应的全部混杂因子 $L$（即达成条件可交换性所需的全部变量）。此外，我们假定正值性成立。接下来的问题是：当 $X$ 维数很高、或包含多个连续变量时，如何在实践中估计这一效应 $\E[Y^{a=1}]-\E[Y^{a=0}]$。

若我们对 $\E[Y^{a=1}]$ 与 $\E[Y^{a=0}]$ 有良好的估计，二者之差就是 $\E[Y^{a=1}]-\E[Y^{a=0}]$ 的良好估计。因此，为简单起见，我们集中讨论 $\E[Y^{a=1}]$ 的估计。

视所选的校正方法不同，变量 $X$ 会以不同方式被使用。用代入式 g 公式（标准化）估计 $\E[Y^{a=1}]$ 时，我们要估计在 $A=1$ 的个体中结局 $Y$ 给定变量 $X$ 的条件均值，记作 $b(X)$；用逆概率加权时，我们要估计处理 $A$ 给定变量 $X$ 的条件概率，记作 $\pi(X)$。我们可以用本书第二部分描述的那类传统参数模型（例如带线性、logistic 或对数连结的广义线性模型）得到估计 $\hat b(x)$ 与 $\hat\pi(x)$，其参数个数远少于样本量。当 $X$ 维数很高时，这类模型必定设定错误。于是 $\hat b(x)$ 与 $\hat\pi(x)$ 都不会是真实 $b(x)$ 与 $\pi(x)$ 的一致估计。

为降低模型设定错误的可能性，我们或许想拟合参数丰富的广义线性模型，其线性预测子为 $\theta^{\mathrm{T}} s(x)=\sum_j \theta_j s_j(x)$，其中 $s(X)$ 是协变量向量 $X$ 的一个维数极高的变换向量。向量 $s(X)$ 通常既含 $X$ 中单个变量的灵活高维变换（如三次样条），也含这些变换的跨变量乘积。具体而言，设我们选用 logit 连结，则
\[
\hat b(x)=\mathrm{expit}\bigl(\hat\theta_b^{\mathrm{T}} s(x)\bigr),\qquad
\hat\pi(x)=\mathrm{expit}\bigl(\hat\theta_\pi^{\mathrm{T}} s(x)\bigr),
\]
其中 $\hat\theta_b$ 与 $\hat\theta_\pi$ 分别是 $b(x)$ 与 $\pi(x)$ 之模型的参数。即便基于这些模型的估计函数 $\hat b(x)$ 与 $\hat\pi(x)$ 是一致的，在有限样本中，误差 $\hat b(x)-b(x)$ 与 $\hat\pi(x)-\pi(x)$ 也会远大于正确设定的低维参数模型的情形。事实上，$s(X)$ 的维数常常可能超过为研究提供数据的个体数 $n$。此时模型拟合将不收敛，不会返回 $\theta$ 的任何估计。\mnote{记住，$X$ 中有些变量甚至可能不是混杂因子，若我们知道它们是哪些，本无需对它们做校正。}

可行的出路有：给线性预测子为 $\theta^{\mathrm{T}} s(X)$ 的参数模型加上 lasso 或岭惩罚（见精要 18.1）后再拟合；采用逐步选择之类的变量选择算法；或用其他预测性机器学习算法（如基于树的算法——随机森林，或神经网络——深度学习）来估计条件期望 $b(X)$ 与 $\pi(X)$。如精要 18.2 所述，深度学习算法所拟合的模型常含成千上万个参数。其他机器学习算法实际上也拟合成千上万个参数。在样本量大、协变量 $X$ 多的多数情形下，机器学习算法在准确预测条件期望方面优于传统参数模型。\mnote{机器学习算法可用交叉验证（见精要 18.2）来优化预测精度。}

然而，仅靠预测性机器学习算法本身，并不足以在高维情境中充分校正混杂。下一节说明，这些算法必须与双重稳健估计量结合使用，并做两处修改：样本拆分与交叉拟合。若我们希望构造有效的 95\% Wald 置信区间——即至少有 95\% 的概率包含所关心的因果参数的区间——这两处修改就是必需的。

\section{双重稳健机器学习估计量}

$\psi=\E[Y^{a=1}]$ 的有效 Wald 区间要求估计量的偏倚远小于其标准误。多数估计量 $\hat\E[Y^{a=1}]$ 的标准误按 $1/\sqrt{n}$ 乘以某常数的方式缩放，其中 $n$ 是样本量。因此，我们要求 $\hat\E[Y^{a=1}]$ 的偏倚远小于 $1/\sqrt{n}$。除偏倚小之外，要得到以 $\hat\E[Y^{a=1}]$ 为中心的有效 Wald 区间，一般还需要它渐近正态，而这通常比小偏倚更容易做到。

用双重稳健估计量比用单一（非双重）稳健估计量更容易做到小偏倚，因为双重稳健估计量的偏倚 $\hat\E[Y^{a=1}]-\E[Y^{a=1}]$ 取决于误差 $\frac{1}{\hat\pi(x)}-\frac{1}{\pi(x)}$ 与 $b(x)-\hat b(x)$ 的乘积，而该乘积可以很小。的确，只要两个误差都远小于 $1/\sqrt[4]{n}$，偏倚就小于 $1/\sqrt{n}$；当函数 $\pi(x)$ 与 $b(x)$ 相当平滑（即有许多阶导数）或非常稀疏（即只依赖向量 $X$ 的少数分量，尽管有多少、是哪些分量未知）时，机器学习估计量常能做到这一点。\mnote{当超总体中存在一定程度的混杂时，覆盖不足会更严重，因为此时 Wald 置信区间并不以因果效应的无偏估计量为中心（见第 10 章）。}

相比之下，$\E[Y^{a=1}]$ 的逆概率加权估计量与代入式 g 公式估计量的偏倚，分别与误差 $\frac{1}{\hat\pi(x)}-\frac{1}{\pi(x)}$ 和 $b(x)-\hat b(x)$ 一样大。若果真如此，逆概率加权估计量与代入式 g 公式估计量一般都无法以 $\E[Y^{a=1}]$ 为中心给出有效的 Wald 置信区间，因为在高维 $X$ 下，已知 $\pi(x)$ 或 $b(x)$ 的任何可能估计量 $\hat\pi(x)$ 或 $\hat b(x)$ 的误差都必定超过 $1/\sqrt{n}$。\mnote{双重稳健估计量的这一性质称为\emph{二阶偏倚}。详见技术注记 13.2。}

但是，若希望构造有效的 95\% Wald 置信区间，上一段所说的双重稳健机器学习估计量就需引入样本拆分与交叉拟合。下面描述这两种程序及其道理。技术注记 18.1 概括估计过程的步骤。

先描述样本拆分。第一，把 $n$ 个个体构成的研究总体随机分成两半：大小为 $n/2$ 的估计样本，以及同样大小的训练样本。第二，把预测性机器学习算法用于训练样本，得到条件期望 $b(x)=\E[Y\given X=x,A=1]$ 与 $\pi(x)=\E[A\given X=x]$ 的估计量 $\hat b(x)$ 与 $\hat\pi(x)$。第三，用来自训练样本的 $\hat b(x)$ 与 $\hat\pi(x)$，在估计样本中计算平均因果效应的双重稳健估计量。至此，我们得到了研究总体随机一半中平均因果效应的双重稳健机器学习估计。\mnote{我们或可把训练样本称为冗余样本，因为我们用它来估计 $b(X)$ 与 $\pi(X)$ 的冗余回归。精要 15.1 回顾了冗余参数的概念。}

要理解样本拆分的必要性，让我们把拆分样本版本与技术注记 13.2 中增广逆概率加权（AIPW）双重稳健估计量的全样本版本作比较。$\E[Y^{a=1}]$ 的拆分样本 AIPW 估计量所用的估计样本，与仅使用训练样本数据的 $\hat b(x)$、$\hat\pi(x)$ 拆分样本估计量在统计上独立。因此，在技术注记 18.2 所述的弱条件下，该估计量渐近正态，其标准误按 $1/\sqrt{n}$ 缩放，并带有前面提到的乘积偏倚。由此可知，若乘积偏倚小于 $1/\sqrt{n}$，以拆分样本估计量为中心的 Wald 区间就是有效的。

相比之下，$\E[Y^{a=1}]$ 的全样本 AIPW 估计量为
\[
\frac{1}{n}\sum_{i=1}^{n}\left[\hat b(X_i)+\frac{A_i}{\hat\pi(X_i)}\bigl\{Y_i-\hat b(X_i)\bigr\}\right],
\]
其中 $\hat b(x)$ 与 $\hat\pi(x)$ 现由施加于全部 $n$ 个个体数据的机器学习算法估计。于是 $\hat b(x)$ 与 $\hat\pi(x)$ 和全样本 AIPW 估计量相关。这种相关若足够强，会以难以预料的方式影响全样本估计量的偏倚、方差与渐近正态性。遗憾的是，相关的大小未知，也无法很好地估计。因此，在高维情境中，拆分样本估计量远比全样本估计量可取。

使用双重稳健拆分样本估计量的唯一困难在于，其方差与标准误对应的是 $n/2$ 的样本量。因此，我们的置信区间会比能用全部 $n$ 个个体时得到的区间更宽。克服这一困难的办法是交叉拟合。\mnote{样本拆分与交叉拟合并非新程序。但把二者与机器学习结合起来的思想，直到晚近才受到重视。}

\renewcommand{\TPname}{技术注记（Technical Point）}
\begin{techpoint}[增广逆概率加权拆分样本与交叉拟合估计量]
$\psi=\E[Y^{a=1}]$ 的增广逆概率加权（AIPW）估计量是技术注记 13.2 所述的双重稳健估计量。下列算法计算 AIPW 拆分样本估计量 $\hat\psi$ 与交叉拟合估计量 $\hat\psi_{\text{cross-fit}}$：

\begin{enumerate}[label=(\roman*)]
\item 把 $n$ 名研究对象随机分成两部分：大小为 $q$ 的估计样本，与大小为 $n_{\mathrm{tr}}=n-q$、且 $q/n\approx 1/2$ 的训练样本。
\item 用机器学习算法从训练样本数据估计 $b(x)=\E[Y\given A=1,X=x]$ 与 $\pi(x)=\pr[A=1\given X=x]$ 的 $\hat b(x)$ 与 $\hat\pi(x)$。
\item 由估计样本中的 $q$ 名对象计算拆分样本 AIPW 估计量
\[
\hat\psi=\frac{1}{q}\sum_{i=1}^{q}\left[\hat b(X_i)+\frac{A_i}{\hat\pi(X_i)}\bigl\{Y_i-\hat b(X_i)\bigr\}\right].
\]
\item 计算交叉拟合估计量
\[
\hat\psi_{\text{cross-fit}}=\bigl(\hat\psi+\tilde\psi\bigr)\big/2,
\]
其中 $\tilde\psi$ 是交换训练样本与估计样本后的 $\hat\psi$。
\end{enumerate}

另一种有限样本表现更佳的交叉拟合估计量按如下步骤计算：(i) 把大小为 $n$ 的样本分成 $M>2$ 个等大的随机样本；(ii) 以样本 $m$ 为估计样本、其余 $M-1$ 个样本为训练样本，计算 $\hat\psi^{(m)}$，$m=1,2,\ldots,M$；(iii) 计算 $\hat\psi_{\text{cross-fit}}=\frac{1}{M}\sum_{m=1}^{M}\hat\psi^{(m)}$。
\end{techpoint}
\renewcommand{\TPname}{技术注记}

下面描述交叉拟合如何挽回样本拆分所损失的统计效率。第一，重复上述程序，但交换研究总体中估计半区与训练半区的角色。也就是说，把原先留作估计的那一半用作新的训练样本，把原先用于训练的那一半用作新的估计样本。然后，用来自新训练样本的 $\hat b(x)$ 与 $\hat\pi(x)$，在新的估计样本中计算平均因果效应的双重稳健估计。至此，我们得到了总体另一半中平均因果效应的双重稳健机器学习估计。

下一步是计算来自总体两半的两个双重稳健估计的平均值。这个平均值就是整个研究总体中效应的双重稳健估计。围绕该估计的 95\% 置信区间可用自助法构造：或者加减 1.96 倍的自助法标准误，或者用自助法估计的 2.5 与 97.5 百分位数作为区间的上下界。

至此完成。通过样本拆分与交叉拟合，我们可以把双重稳健估计与机器学习结合起来，得到具有已知统计性质、且使用了全部可得数据的因果效应估计。一个活跃的研究方向是：发展程序来检测双重稳健拆分样本估计量的偏倚是否达到或超过标准误的量级，若是，则在估计样本中获得偏倚更小的估计，而不必在训练样本中重做机器学习部分。\mnote{Lin 等 (2020) 基于高阶影响函数构造了估计量，其偏倚小于双重稳健交叉拟合估计量，而方差没有显著增大。}

\begin{techpoint}[拆分样本与交叉拟合估计量的统计性质]
给定训练样本数据 $Tr$，$\hat b(x)$ 与 $\hat\pi(x)$ 是固定函数。因此 $\hat\psi$ 是独立同分布随机变量之和，由中心极限定理，它在给定 $Tr$ 的条件下渐近正态，标准误 $\mathrm{se}\bigl(\hat\psi\bigr)$ 与 $n^{-1/2}$ 成正比。$\hat\psi$ 的精确条件偏倚为
\[
\E\bigl[\hat\psi-\psi\given Tr\bigr]
=\E\Bigl[\pi(X_i)\Bigl(\tfrac{1}{\hat\pi(X_i)}-\tfrac{1}{\pi(X_i)}\Bigr)
\bigl\{b(X)-\hat b(X)\bigr\}\given Tr\Bigr].
\]

要刻画 $\hat\psi$ 与 $\hat\psi_{\text{cross-fit}}$ 的无条件统计性质，必须考虑到 $\E\bigl[\hat\psi-\psi\given Tr\bigr]$ 通过 $\hat b$ 与 $\hat\pi$ 对训练样本数据的依赖而是随机的。若 (i) $\hat b(x)$ 与 $\hat\pi(x)$ 分别是真实 $b(x)$ 与 $\pi(x)$ 的一致估计（在均方意义下），且 (ii) $\E\bigl[\hat\psi-\psi\given Tr\bigr]/\mathrm{se}\bigl(\hat\psi\bigr)$ 依概率收敛到 0，则 $\hat\psi$ 与 $\hat\psi_{\text{cross-fit}}$ 渐近正态且无偏。

于是，当 (i) 与 (ii) 成立时，95\% Wald 置信区间 $\hat\psi\pm 1.96\times\hat{\mathrm{se}}\bigl(\hat\psi\bigr)$ 与 $\hat\psi_{\text{cross-fit}}\pm 1.96\times\hat{\mathrm{se}}\bigl(\hat\psi_{\text{cross-fit}}\bigr)$ 是有效的，而且实际上是标定的。此处 $\hat{\mathrm{se}}\bigl(\hat\psi\bigr)$ 与 $\hat{\mathrm{se}}\bigl(\hat\psi_{\text{cross-fit}}\bigr)$ 可用自助法计算。进一步，$n^{1/2}\hat{\mathrm{se}}\bigl(\hat\psi_{\text{cross-fit}}\bigr)$ 是半参数有效的，其标准误为
\[
\Bigl[\Var\Bigl\{b(X)+\frac{A}{\pi(X)}\bigl[Y-b(X)\bigr]\Bigr\}\Bigr]^{1/2},
\]
它比 $\hat\psi$ 的标准误小 $1/\sqrt{2}$ 倍。注意，若 $\frac{1}{\hat\pi(x)}-\frac{1}{\pi(x)}$ 的收敛速度为 $n^{-\alpha}$，$b(x)-\hat b(x)$ 的收敛速度为 $n^{-\epsilon}$，则当 $\alpha+\epsilon>1/2$ 时，偏倚 $\E\bigl[\hat\psi-\psi\given Tr\bigr]$ 为 $o\bigl(n^{-1/2}\bigr)$。因此，若 $\hat b(x)$ 的收敛速度慢于 $n^{-1/4}$，只要 $\hat\pi(x)$ 的收敛速度足够快，偏倚仍可为 $o\bigl(n^{-1/2}\bigr)$。把 $\hat b(x)$ 与 $\hat\pi(x)$ 的角色互换，结论同样成立。
\end{techpoint}

\section{变量选择是一个难题}

上一节概述的方法，推翻了那种广为流传的看法：任何数据自适应地挑选校正变量的程序，都必然导致错误的置信区间。如我们所见，把因果推断方法与机器学习算法结合起来挑选混杂因子，在一定条件下可以得到正确的统计推断。然而，双重稳健机器学习并未解决我们的全部问题，理由至少有三（此外还有上一节所述的那一条）。

第一，在许多应用中，可得的专业知识可能不足以识别全部重要混杂因子，也不足以排除那些会诱发或放大偏倚的变量。因此，无法保证双重稳健机器学习估计量的偏倚一定很小。

第二，在含时变处理的高维情境中，双重稳健估计量的实现一直很困难——与机器学习结合时计算代价高昂。因果生存分析尤其如此。结果，已发表的复杂纵向数据因果推断例子，多数使用单一稳健估计量，也就是本书第三部分着重介绍的那些。不过，本章概述的方法在某些领域正迅速成为常规做法。

第三，双重稳健机器学习可以给出一个方差，它等于真实条件期望 $b(X)$ 与 $\pi(X)$ 已知时会得到的方差。然而，无法保证这样的方差小到足以做有意义的因果推断。

设想我们按上一节所述得到了因果效应的双重稳健机器学习估计，结果却发现其（正确的）方差太大而无甚用处。即便我们用很小的乘积偏倚估计了倾向性评分与结局回归，只要 $X$ 中某些协变量与处理 $A$ 强相关，这种情况仍会发生。此时，对具有特定 $X$ 取值的个体，处理概率 $\pi(X)$ 可能接近 0 或接近 1。结果，效应估计的方差极大，其 95\% 置信区间也就极宽（但往往是正确的）。

既然我们不喜欢极宽的 95\% 置信区间——即便它们是正确的——就可能忍不住把 $X$ 中造成“问题”的变量扔掉，再重做数据分析。若真那样做，我们就根本改变了游戏规则。用数据去剔除 $X$ 中与处理相关、而与结局关系不那么大的协变量，就再也无法保证效应估计周围的 95\% 置信区间是有效的。为消除偏倚而纳入全部潜在混杂因子，与为减小方差而排除某些变量，这二者之间的张力很难化解。\mnote{这一结果引出一个令人困惑的哲学问题：若我们用数据剔除了（譬如说）5 个使方差过大的变量，置信区间就失效，那么，倘若我们碰巧拿到的数据集本就不含这 5 个变量，置信区间为何又该有效？既然我们总是在某些潜在混杂因子未被记录的数据集上工作，究竟该如何解释任何观察性分析中的置信区间？}

鉴于以上种种，要为变量选择制定一套清晰的一般准则或许并不可能。事实上，围绕这些议题的方法论研究如此之多，本章不可能给出规定性的结论。如第 13.5 节所述，对因果推断最好的科学建议或许是开展多重敏感性分析：实施若干种分析方法，检视所得的效应估计。若各种效应估计彼此相容，我们对结果会更有信心；若各种效应估计并不相容，作为研究者，我们的任务就是设法弄清原因。
